%% 04-theorem-a
\section{A No-Go Theorem for Product Encodings}
\label{sec:theorem-a}

\begin{theorem}[Product no-go]
  \label{thm:product-nogo}
  Let $\Phi$ be a block-product feature map on $\Zpn$ (\Cref{def:block-product-map})
  whose fidelity kernel is ultrametric with profile $f$, and set
  $v^{*} = \min\{\, v : f(v) > 0 \,\}$. Let $B^{*}$ be the block containing level
  $v^{*}+1$. Then for every block $B$ all of whose levels exceed $\max B^{*}$, and all
  digit tuples $a, b$,
  \begin{equation}
    \label{eq:trivial-block}
    \abs{\braket{\psi_B(a)}{\psi_B(b)}} = 1 .
  \end{equation}
  Consequently $f$ takes at most the three values $\{0, f(v^{*}), 1\}$, and $\Phi$
  resolves the tree only to depth $\max B^{*}$. In particular, for a product map with
  $n \ge 2$, no strictly monotone profile is realisable.
\end{theorem}

\begin{proof}
  The argument uses exactly one property of the encoding: fidelity is multiplicative
  across tensor factors, \cref{eq:multiplicative}. We give it for $v^{*} = 0$; the
  general case shifts every index by $v^{*}$, since all pairs meeting above depth
  $v^{*}$ have kernel value $0$ and impose no constraint.

  Let $B$ be a block whose levels all exceed $\max B^{*}$, and fix digit tuples
  $a \ne b$ for it. Choose leaves $x, y$ that differ in level $1$ and agree at every
  other level, with $x_B = a$. Then $\lcadepth(x,y) = 0$, and since every factor other
  than $B^{*}$ contributes $g(x, x) = 1$,
  \begin{equation}
    \label{eq:base-value}
    f(0) = K(x, y) = \fid{\psi_{B^{*}}(x_{B^{*}})}{\psi_{B^{*}}(y_{B^{*}})} .
  \end{equation}

  Now let $y'$ agree with $y$ everywhere except on $B$, where $y'_B = b$. Because $x$
  and $y'$ still differ at level $1$, and $B$ does not contain level $1$, we still have
  $\lcadepth(x, y') = 0$ and therefore $K(x, y') = f(0)$. But multiplicativity gives
  \begin{equation*}
    K(x, y')
    = \fid{\psi_{B^{*}}(x_{B^{*}})}{\psi_{B^{*}}(y_{B^{*}})} \cdot
      \fid{\psi_{B}(a)}{\psi_{B}(b)}
    = f(0) \cdot \fid{\psi_B(a)}{\psi_B(b)} .
  \end{equation*}
  Since $f(0) > 0$ by the choice $v^{*} = 0$, this forces
  $\fid{\psi_B(a)}{\psi_B(b)} = 1$, which is \cref{eq:trivial-block}: the states of
  block $B$ coincide up to a global phase, and $B$ carries no information about the
  data.

  Every block beyond $B^{*}$ being trivial, $K$ depends only on the digits in blocks up
  to $B^{*}$. Two leaves agreeing down to depth $\max B^{*}$ therefore have
  $K(x,y) = 1$, so $f(v) = 1$ for all $v > \max B^{*}$ --- the tree is resolved no
  deeper. Level constancy forces the common value $f(0)$ on all pairs meeting at the
  root, so together with the vanishing values below $v^{*}$ the profile takes at most
  the three values $\{0, f(v^{*}), 1\}$. For a product map every block is a single
  level, $\max B^{*} = v^{*}+1$, and a strictly monotone profile would need $n+1$
  distinct values, which exceeds three whenever $n \ge 2$.
\end{proof}

The proof turns on multiplicativity alone. It is therefore indifferent to the local
dimensions $d_i$, to whether the local states are real or complex, and to how they were
chosen --- which is what makes it a statement about a class of encodings rather than
about particular instances.

\begin{corollary}
  \label{cor:standard-encodings}
  On $\Zpn$ with $n \ge 2$, none of the following induces a strictly monotone
  ultrametric kernel: angle encoding with one qubit per digit; computational-basis
  encoding; and any first-order Pauli-$Z$ feature map, whose phase is a sum of
  single-digit terms and which is therefore a product map.
\end{corollary}

Basis encoding is the extreme case and deserves a sentence, because it is the one that
can be mistaken for a success. It attains $v^{*} = n$: its kernel is the identity, its
induced distance is the discrete metric, and it satisfies \cref{eq:strong-triangle}
with zero violations while resolving nothing. It is a corner of
\Cref{thm:product-nogo}, not a counterexample to it, and it is the reason
\Cref{subsec:diagnostics} insists on reporting resolution depth alongside violations.

\begin{remark}[What this theorem does not cover]
  \label{rem:not-covered}
  A ZZ or IQP feature map applies all-pairs couplings between qubits, so its state does
  not factorise across digits and \cref{eq:multiplicative} fails. \Cref{thm:product-nogo}
  says nothing about it. We do not patch the gap by weakening the claim; we close it in
  \Cref{sec:theorem-b} with an argument that ignores the structure of the encoding
  entirely.
\end{remark}

\paragraph{Numerical verification.}
Randomly drawn product maps are never level-constant for
$(p, n) \in \{(2,3), (3,2), (2,4), (5,2)\}$ across five seeds
(\texttt{test\_theorem\_a\_random\_product\_maps\_are\_never\_strictly\_monotone}). A
product map built to satisfy the conclusion --- one non-trivial level, the rest trivial
--- yields exactly two distinct kernel values, as predicted
(\texttt{test\_theorem\_a\_ultrametric\_product\_map\_resolves\_only\_one\_level}), and
switching on a second non-trivial level destroys level constancy, which is the
contradiction the proof derives
(\texttt{test\_theorem\_a\_a\_nontrivial\_deep\_factor\_destroys\_ultrametricity}).
